\chapter{PROBLEM DEFINITION AND RESEARCH QUESTIONS}

% ================================================
% 3.1 Introduction
% ================================================
\section{Introduction}

This chapter formalizes the research problem addressed by this thesis. Based on empirical analysis of the NYC Yellow Taxi dataset, we identify critical gaps in existing streaming DQ frameworks and derive research questions that guide the rest of this work.

% ================================================
% 3.2 Case Study: NYC Yellow Taxi Dataset
% ================================================
A detailed exploratory analysis of the NYC Yellow Taxi dataset is provided in Appendix A. The analysis covers:
\begin{itemize}
    \item \textbf{Dataset schema and structure} (Section A.1): 19 raw columns plus 4 derived features.
    \item \textbf{Temporal coverage and volume patterns} (Section A.1): 3,076,903 records from July 2024, with validation across 26 months (January 2024 -- February 2026).
    \item \textbf{Statistical distributions} (Section A.1): Right-skewed fare distribution (mode at \$8--\$12 for standard trips, \$50--\$60 for airport trips).
    \item \textbf{Data quality observations} (Section A.1): 9.1\% NULL rate, 1.8\% negative fares, 1.5\% zero distances.
    \item \textbf{Schema validation results} (Section A.2): 278,989 records (9.07\%) fail completeness checks.
    \item \textbf{Business rules results} (Section A.3): 110,941 records (3.87\%) fail validity checks.
    \item \textbf{Multivariate anomaly patterns} (Section A.4): Four fraud use cases undetectable by single-feature rules.
\end{itemize}

These findings establish the empirical basis for the research gaps and research questions stated in this chapter.

% ================================================
% 3.6 Research Gaps
% ================================================
\section{Research Gaps}

Based on the case study analysis (Appendix A), three research gaps are identified:

\subsection{Gap 1: Context Collapse in Heterogeneous Data}

The NYC Taxi dataset exhibits extreme heterogeneity across multiple dimensions (Appendix A, Tables A.5--A.7). Airport trips (RatecodeID=2,3) have average fares three to four times higher than standard trips, and time-of-day patterns shift the distribution of special ratecodes significantly. A single global threshold calibrated on standard trips produces catastrophically high false positive rates on airport trips, while remaining too permissive for standard trips. Existing open-source DQ frameworks (Great Expectations, AWS Deequ, Soda Core) provide no mechanism for assigning different thresholds to different data contexts. This gap manifests as the inability to distinguish "legitimately high-value" trips from fraudulent ones without manual per-context rule specification.

\subsection{Gap 2: Sequential Pipelines Underutilize Heterogeneous Validation Complexity}

Existing streaming DQ frameworks process records through validation stages sequentially: each record must complete Schema Validation before Business Rules, and Business Rules before ML scoring. This architectural choice means that all records pay the full cost of the slowest validation stage, regardless of whether earlier stages could make a determination. For taxi data, initial validation estimates that approximately 10--15\% of records fail Business Rules, routing these records through ML scoring wastes computational resources on records already known to be violations. No existing framework parallelizes heterogeneous validation stages to exploit early exit opportunities.

\subsection{Gap 3: Single-Strategy Drift Adaptation Fails to Differentiate Drift Severity}

When concept drift occurs in production DQ systems, existing frameworks apply a uniform response: full model retraining on the entire historical dataset. This approach is computationally expensive and introduces significant downtime during which the stale model continues to produce unreliable results. Our case study shows that NYC Taxi data experiences multiple drift types: gradual seasonal shifts (months), sudden event-driven spikes (holidays, weather), and recurring cyclical patterns (weekday vs. weekend). Treating all drift types identically wastes resources on minor shifts that could be resolved by threshold adjustment, while potentially under-responding to severe shifts that require full retraining. No existing framework classifies drift type to select the appropriate adaptation strategy.

\subsection{Qualitative Comparison with Existing DQ Frameworks}

To establish the novelty of CA-DQStream, we compare its capabilities against existing data quality frameworks across four key dimensions: ML-based detection, streaming architecture, context-aware thresholds, and drift adaptation. This qualitative comparison demonstrates that CA-DQStream is the only framework combining all four capabilities.

\begin{longtable}{|m{2.8cm}|m{2cm}|m{1.6cm}|m{1.6cm}|m{1.6cm}|m{1.8cm}|}
\hline
\textbf{Framework} & \textbf{Type} & \textbf{ML Detection} & \textbf{Streaming} & \textbf{Context-Aware} & \textbf{Drift Adapt} \\ \hline
Great Expectations & Batch & \ding{55} & \ding{55} & \ding{55} & \ding{55} \\ \hline
Soda Core & Batch & \ding{55} & \ding{55} & \ding{55} & \ding{55} \\ \hline
AWS Deequ & Batch & \ding{55} & \ding{55} & \ding{55} & \ding{55} \\ \hline
Stream DaQ & Streaming & \ding{55} & \checkmark & \ding{55} & \ding{55} \\ \hline
MemStream & Streaming & \checkmark & \checkmark & \ding{55} & \checkmark \\ \hline
MStream & Streaming & \checkmark & \checkmark & \ding{55} & \checkmark \\ \hline
\textbf{CA-DQStream (Ours)} & \textbf{Streaming} & \textbf{\checkmark} & \textbf{\checkmark} & \textbf{\checkmark} & \textbf{\checkmark} \\ \hline
\caption{Qualitative comparison of CA-DQStream against existing DQ frameworks.}
\label{tab:dq_framework_comparison}
\end{longtable}

\medskip
\noindent Great Expectations, Soda Core, and AWS Deequ are batch-oriented DQ frameworks; they are fundamentally incomparable to streaming systems in throughput and latency. They are included for completeness only. The primary streaming DQ comparison focuses on MemStream, MStream, and Stream DaQ, which share similar problem settings. Unlike these prior systems, CA-DQStream combines (1) multi-dimensional context-aware thresholding, (2) fork-join parallel validation, (3) multi-strategy drift adaptation, and (4) hybrid clustering-based anomaly detection in a unified framework. No existing streaming DQ system achieves all four capabilities simultaneously.\footnote{The streaming DQ comparison focuses on systems that operate in the same data paradigm: MemStream \citep{bhatia2022memstream}, MStream \citep{siddiqui2019mstream}, and Stream DaQ \citep{papastergios2025stream}. MemStream and MStream provide streaming anomaly detection but lack context-aware thresholds. Stream DaQ provides streaming DQ monitoring but lacks ML-based anomaly detection and drift adaptation.}

\subsection{Ground Truth and Evaluation Methodology}

A critical methodological consideration in data quality research is the absence of labeled ground truth for real-world anomalies. We adopt the standard practice of the data quality community: evaluating on \textbf{synthetic injected anomalies} with controlled properties.

Using synthetic injected anomalies is \textbf{standard practice} in data quality and data cleaning research, explicitly used in top-tier publications:

\begin{itemize}
    \item \textbf{Information Systems, Q1 2022} \citep{DCRepair2022}: Authors use an ultra-clean baseline and intentionally inject DC violations using data pollution tools to ensure fair comparison across baseline methods. They explicitly state: \emph{``All baseline methods use injected errors to ensure fair comparison.''}

    \item \textbf{Information Systems, Q1 2024} \citep{DCRepair2024}: Authors justify synthetic injection with two key reasons: (1) real-world errors rarely come with ground truth labels, making precision/recall measurement impossible; (2) all baselines must use identical injected errors to guarantee fair comparison.

    \item \textbf{Environmental Modelling \& Software, 2013} \citep{bayesianQC2013}: Authors generate synthetic data errors (missed tip, false tip) on weather data streams, acknowledging: \emph{``It is difficult to know the true classification of natural observed data to calculate false positive/negative rates.''}

    \item \textbf{EDBT, 2025} \citep{icewafl2025}: The entire paper introduces \emph{Icewafl}, a configurable data stream polluter. Authors assert that injection tools enable systematic evaluation of algorithms against diverse data error types.

    \item \textbf{Ada-Context, 2025} \citep{adacontext2025}: Authors manipulate 5\% of test data by injecting inaccurate, invalid, and incomplete errors following ISO 25012 standard, specifically to measure cleaning algorithm effectiveness.

    \item \textbf{VLDB, 2024} \citep{dafdiscover2024}: Authors clone real data and use the Bart tool to inject errors (adding ``*'' characters) at probabilistic rates to simulate real data quality issues.
\end{itemize}

\textbf{Justification for CA-DQStream:} Ground truth is unavailable for real-world taxi fraud data. NYC TLC does not publish labeled fraud cases. Using synthetic injected anomalies allows us to: (1) precisely control anomaly difficulty (Easy/Medium/Hard), (2) ensure fair comparison across all baseline methods, and (3) measure precision, recall, and FPR with known denominators. This approach is consistent with established practice in Information Systems Q1 venues, EDBT, and VLDB.

% ================================================
% 3.7 Research Questions
% ================================================
\section{Research Questions}

This thesis addresses five research questions:

\subsection{RQ1: Multi-Layer Validation}
\textit{``Does a multi-layer architecture (Schema + Rules + ML + Drift) improve data quality coverage compared to single-layer validation?''}

\begin{itemize}
    \item \textbf{RQ1a:} How many violations does each processing layer catch, and what is the additive coverage of the multi-layer architecture?
    \item \textbf{RQ1b:} What proportion of anomalies are unique to each layer (i.e., caught by only one layer)?
    \item \textbf{RQ1c:} Does removing any layer significantly degrade FPR or coverage?
\end{itemize}

\subsection{RQ2: Hybrid K-Means + Online Autoencoder Model}
\textit{``Does the Context-Aware Online Autoencoder (21D ratio features + per-cluster thresholds) outperform simpler baseline approaches and established streaming algorithms on multivariate taxi trip anomalies?''}

\begin{itemize}
    \item \textbf{RQ2a:} What is the incremental value of ratio features over raw features?
    \item \textbf{RQ2b:} What is the incremental value of K-Means clustering over global thresholds?
    \item \textbf{RQ2c:} Does the full Hybrid model achieve production targets (FPR $<$ 5\%, Recall $\geq$ 75\%)?
\end{itemize}

\subsection{RQ3: IEC Multi-Strategy Adaptation}
\textit{``Does the Intelligent Evolution Controller's multi-strategy adaptation (Continuous Evolution, Switching, METER, Spatial Tracking) maintain model accuracy when data distributions drift over time?''}

\begin{itemize}
    \item \textbf{RQ3a:} How quickly does ADWIN detect drift events on quality meta-streams?
    \item \textbf{RQ3b:} Does the multi-strategy controller select appropriate strategies based on drift characteristics?
    \item \textbf{RQ3c:} Can lightweight strategies (Continuous Evolution, Switching) resolve most drift events without full retraining?
\end{itemize}

\subsection{RQ4: Rendezvous Pipeline Architecture}
\textit{``Does the Rendezvous (fork-join) pipeline architecture reduce latency and increase throughput compared to traditional linear (sequential) pipelines?''}

\begin{itemize}
    \item \textbf{RQ4a:} Does the Rendezvous architecture achieve lower average latency through early exit optimization compared to linear pipelines?
    \item \textbf{RQ4b:} Does the early exit optimization significantly improve throughput?
    \item \textbf{RQ4c:} Does the system degrade gracefully when the ML branch fails?
\end{itemize}

\subsection{RQ5: 4D Context-Aware Thresholding}
\textit{``Does the 4D context-aware threshold matrix (trip\_type $\times$ time\_window $\times$ day\_type $\times$ neighborhood) reduce false positive rate compared to global thresholds?''}

\begin{itemize}
    \item \textbf{RQ5a:} What is the FPR of 4D context-aware thresholds vs single global threshold?
    \item \textbf{RQ5b:} Does the spatial dimension (neighborhood\_id) reduce FPR?
    \item \textbf{RQ5c:} Does the temporal dimension (time\_window $\times$ day\_type) reduce FPR?
    \item \textbf{RQ5d:} How many context cells require fallback to global thresholds due to insufficient training data?
\end{itemize}

% ================================================
% 3.8 Research Hypotheses
% ================================================
\section{Research Hypotheses}

Based on the case study analysis, we state the following hypotheses:

\begin{enumerate}
    \item \textbf{H1:} At least 40\% of total violations detected by the four-stage pipeline are unique to Layer 3 (Online Autoencoder) --- undetectable by Layers 1--2 (Schema Validation + Business Rules) alone. These represent multivariate anomalies that require understanding inter-feature correlations, such as meter tampering, short-trip fraud, and payment inflation patterns.
    \item \textbf{H2:} The combined four-stage pipeline achieves FPR below 5\% (specifically 4.2\% at month 1, rising to 4.9\% at month 26 with IEC recalibration) compared to 38.7\% weighted average for global thresholds (up to 68.9\% for Newark airport trips), and F1 $\geq$ 0.828 on Easy-difficulty anomalies on the NYC Taxi dataset.
    \item \textbf{H3:} The IEC detects drift events within 24 hours (1440 aggregated 1-minute windows) on average, and selects appropriate adaptation strategies based on drift characteristics, with 93\% of drift events resolved by lightweight incremental updates (Continuous Evolution and Switching) without full retraining.
    \item \textbf{H4:} The Rendezvous pipeline architecture reduces average end-to-end latency through early exit optimization and improves overall system throughput through parallel computation, compared to linear sequential pipelines.
    \item \textbf{H5:} 4D context-aware thresholds achieve FPR below 5\% compared to 38.7\% weighted average for global thresholds (up to 68.9\% for Newark airport trips, at least 5$\times$ weighted improvement). The weighted average of 38.7\% is computed as the average of per-category FPRs (Standard 12.4\%, JFK 51.2\%, Newark 68.9\%, Negotiated 42.1\%); the 9.2$\times$ raw ratio (38.7\% $\to$ 4.2\% per-record) reflects the actual reduction for the majority Standard trip type.
    \item \textbf{H6:} The CA-DQStream framework's Flink keyed state management (layer\_stats, cluster\_params, drift\_flags) enables stateful, fault-tolerant processing that preserves pipeline accuracy across checkpoint/restart cycles and supports exactly-once semantics for downstream consumers.
\end{enumerate}

% ================================================
% 3.9 Chapter Summary
% ================================================
\section{Chapter Summary}

This chapter formalized the research problem based on the NYC Taxi dataset case study. The case study (detailed in Appendix A) identified five categories of data quality problems: completeness (9.1\% NULL rate), validity ($\sim$3.4\% hard rule violations), multivariate anomalies ($\sim$2.9\% via Online Autoencoder), uniqueness (duplicate records), and concept drift (NULL rate increasing over 26 months). Three research gaps were identified: lack of context-aware thresholding, absence of self-monitoring, and batch-oriented frameworks. Five research questions and six hypotheses guide the evaluation in Chapter 6, with RQ5 and H6 addressing the core contribution of 4D context-aware thresholding. The next chapter details the CA-DQStream framework design.
